\documentclass[11pt,a4paper]{article}
% preamble.tex — estilo común de todos los apuntes Froth.
% Cada apunte hace % preamble.tex — estilo común de todos los apuntes Froth.
% Cada apunte hace % preamble.tex — estilo común de todos los apuntes Froth.
% Cada apunte hace \input{preamble} justo después de \documentclass.
\usepackage[margin=2.4cm]{geometry}
\usepackage{xcolor}
\definecolor{litblue}{RGB}{31,79,121}
\definecolor{litgreen}{RGB}{27,94,32}
\definecolor{litorange}{RGB}{230,81,0}
\usepackage{parskip}
\usepackage{titlesec}
\titleformat{\section}{\Large\bfseries\color{litblue}}{\thesection}{0.6em}{}
\titleformat{\subsection}{\large\bfseries\color{litblue}}{\thesubsection}{0.6em}{}
\usepackage{tcolorbox}
\usepackage{fancyhdr}
\pagestyle{fancy}
\fancyhf{}
\fancyhead[L]{\small\color{litblue}Froth — Apuntes de ML}
\fancyhead[R]{\small\thepage}
\renewcommand{\headrulewidth}{0.4pt}
\usepackage[colorlinks=true,linkcolor=litblue,urlcolor=litblue]{hyperref}

% Cabecera de cada apunte: \cabecera{numero}{titulo}
\newcommand{\cabecera}[2]{%
  \begin{tcolorbox}[colback=litblue,coltext=white,colframe=litblue,arc=2mm]
    {\Large\bfseries Apunte #1 — #2}\\[3pt]
    {\small Proyecto Froth · aprender Machine Learning construyendo}
  \end{tcolorbox}\vspace{0.8em}}

% Cajas temáticas
\newtcolorbox{concepto}[1]{colback=blue!4,colframe=litblue,title={Concepto: #1},arc=1.5mm}
\newtcolorbox{practica}{colback=green!5,colframe=litgreen,title={Pruébalo tú (teclado en mano)},arc=1.5mm}
\newtcolorbox{ojo}{colback=orange!8,colframe=litorange,title={Ojo / error típico},arc=1.5mm}
\newtcolorbox{resumen}{colback=gray!8,colframe=gray!60,title={En una frase},arc=1.5mm}

\newcommand{\codigo}[1]{\texttt{#1}}
 justo después de \documentclass.
\usepackage[margin=2.4cm]{geometry}
\usepackage{xcolor}
\definecolor{litblue}{RGB}{31,79,121}
\definecolor{litgreen}{RGB}{27,94,32}
\definecolor{litorange}{RGB}{230,81,0}
\usepackage{parskip}
\usepackage{titlesec}
\titleformat{\section}{\Large\bfseries\color{litblue}}{\thesection}{0.6em}{}
\titleformat{\subsection}{\large\bfseries\color{litblue}}{\thesubsection}{0.6em}{}
\usepackage{tcolorbox}
\usepackage{fancyhdr}
\pagestyle{fancy}
\fancyhf{}
\fancyhead[L]{\small\color{litblue}Froth — Apuntes de ML}
\fancyhead[R]{\small\thepage}
\renewcommand{\headrulewidth}{0.4pt}
\usepackage[colorlinks=true,linkcolor=litblue,urlcolor=litblue]{hyperref}

% Cabecera de cada apunte: \cabecera{numero}{titulo}
\newcommand{\cabecera}[2]{%
  \begin{tcolorbox}[colback=litblue,coltext=white,colframe=litblue,arc=2mm]
    {\Large\bfseries Apunte #1 — #2}\\[3pt]
    {\small Proyecto Froth · aprender Machine Learning construyendo}
  \end{tcolorbox}\vspace{0.8em}}

% Cajas temáticas
\newtcolorbox{concepto}[1]{colback=blue!4,colframe=litblue,title={Concepto: #1},arc=1.5mm}
\newtcolorbox{practica}{colback=green!5,colframe=litgreen,title={Pruébalo tú (teclado en mano)},arc=1.5mm}
\newtcolorbox{ojo}{colback=orange!8,colframe=litorange,title={Ojo / error típico},arc=1.5mm}
\newtcolorbox{resumen}{colback=gray!8,colframe=gray!60,title={En una frase},arc=1.5mm}

\newcommand{\codigo}[1]{\texttt{#1}}
 justo después de \documentclass.
\usepackage[margin=2.4cm]{geometry}
\usepackage{xcolor}
\definecolor{litblue}{RGB}{31,79,121}
\definecolor{litgreen}{RGB}{27,94,32}
\definecolor{litorange}{RGB}{230,81,0}
\usepackage{parskip}
\usepackage{titlesec}
\titleformat{\section}{\Large\bfseries\color{litblue}}{\thesection}{0.6em}{}
\titleformat{\subsection}{\large\bfseries\color{litblue}}{\thesubsection}{0.6em}{}
\usepackage{tcolorbox}
\usepackage{fancyhdr}
\pagestyle{fancy}
\fancyhf{}
\fancyhead[L]{\small\color{litblue}Froth — Apuntes de ML}
\fancyhead[R]{\small\thepage}
\renewcommand{\headrulewidth}{0.4pt}
\usepackage[colorlinks=true,linkcolor=litblue,urlcolor=litblue]{hyperref}

% Cabecera de cada apunte: \cabecera{numero}{titulo}
\newcommand{\cabecera}[2]{%
  \begin{tcolorbox}[colback=litblue,coltext=white,colframe=litblue,arc=2mm]
    {\Large\bfseries Apunte #1 — #2}\\[3pt]
    {\small Proyecto Froth · aprender Machine Learning construyendo}
  \end{tcolorbox}\vspace{0.8em}}

% Cajas temáticas
\newtcolorbox{concepto}[1]{colback=blue!4,colframe=litblue,title={Concepto: #1},arc=1.5mm}
\newtcolorbox{practica}{colback=green!5,colframe=litgreen,title={Pruébalo tú (teclado en mano)},arc=1.5mm}
\newtcolorbox{ojo}{colback=orange!8,colframe=litorange,title={Ojo / error típico},arc=1.5mm}
\newtcolorbox{resumen}{colback=gray!8,colframe=gray!60,title={En una frase},arc=1.5mm}

\newcommand{\codigo}[1]{\texttt{#1}}

\begin{document}
\cabecera{04}{Tablas de datos: pandas y el archivo parquet}

\section{Qué conseguimos hoy}

El primer botín del proyecto: \textbf{126 papers reales} sobre tu topic, guardados en
\codigo{data/processed/papers.parquet}. Este apunte explica qué es esa ``tabla'' y por qué
se guarda así.

\begin{concepto}{DataFrame (pandas)}
\codigo{pandas} es LA librería de tablas en Python, y su objeto estrella es el
\textbf{DataFrame}: una tabla con filas y columnas, como una hoja de Excel pero manejada
con código. La nuestra tiene una fila por paper y estas columnas:

\begin{center}
\begin{tabular}{ll}
\codigo{id}, \codigo{doi} & identificadores únicos del paper \\
\codigo{title}, \codigo{year} & título y año \\
\codigo{citations} & cuántas veces lo han citado \\
\codigo{authors} & autores (texto ``A; B; C'') \\
\codigo{abstract} & el resumen — \textbf{la materia prima del ML} \\
\codigo{query} & con qué búsqueda apareció (para depurar) \\
\end{tabular}
\end{center}

Con un DataFrame haces en una línea cosas que en Excel serían dolorosas:
\codigo{df.drop\_duplicates(subset="id")} quita papers repetidos;
\codigo{df[df["abstract"].str.len() > 20]} filtra los que no tienen resumen;
\codigo{df.sort\_values("citations")} ordena por citas.
\end{concepto}

\section{La limpieza importa más de lo que parece}

\codigo{pull.py} no guarda lo que llega tal cual. Hace tres limpiezas:
\begin{enumerate}
  \item \textbf{Deduplicar}: el mismo paper puede aparecer en varias búsquedas
        (``lepidolite flotation'' y ``lithium mica flotation'' traen solapados).
  \item \textbf{Filtrar sin-abstract}: sin resumen no hay significado que vectorizar.
  \item \textbf{Reconstruir el abstract} (ver Apunte 02).
\end{enumerate}

\begin{ojo}
En nuestros 126 papers se colaron algunos \textbf{off-topic} (p.\ ej.\ uno de
telecomunicaciones). ¿Por qué? Colisión de siglas: buscamos ``LCT'' (un tipo de pegmatita)
y también existe en otros campos. No lo arreglamos a mano: en las Fases 2--3 el propio ML
los apartará, porque sus \emph{vectores} quedarán lejos de los de flotación. Los datos
reales siempre traen ruido; el sistema debe tolerarlo.
\end{ojo}

\begin{concepto}{parquet vs.\ CSV}
Podíamos guardar la tabla como \codigo{.csv} (texto plano separado por comas). Elegimos
\textbf{parquet} porque:
\begin{itemize}
  \item \textbf{Comprime}: ocupa varias veces menos.
  \item \textbf{Guarda por columnas}: leer solo \codigo{year} y \codigo{citations} de
        100\,000 papers no obliga a leer los abstracts.
  \item \textbf{Conserva los tipos}: en CSV todo vuelve como texto y hay que re-adivinar
        qué era número; parquet recuerda que \codigo{year} es un entero.
\end{itemize}
Es el formato estándar en data science moderno. La librería \codigo{pyarrow} es quien
sabe leerlo/escribirlo; pandas la usa por debajo con
\codigo{df.to\_parquet()} y \codigo{pd.read\_parquet()}.
\end{concepto}

\begin{practica}
Mira tu botín con tus propias manos. En PowerShell, dentro de la carpeta del proyecto,
abre un Python interactivo con \codigo{.venv\textbackslash Scripts\textbackslash python.exe}
y escribe línea a línea:
\begin{enumerate}
  \item \codigo{import pandas as pd}
  \item \codigo{df = pd.read\_parquet("data/processed/papers.parquet")}
  \item \codigo{len(df)} \quad (¿cuántos papers?)
  \item \codigo{df.columns} \quad (las columnas)
  \item \codigo{df.sort\_values("citations", ascending=False).head(3)[["title","year"]]}
  \item \codigo{exit()} para salir.
\end{enumerate}
Acabas de hacer tu primer análisis de datos en pandas.
\end{practica}

\begin{resumen}
Un DataFrame es una tabla programable; limpiamos duplicados y sin-abstract; y parquet es
el formato comprimido, por columnas y con tipos que usa el data science serio.
\end{resumen}

\end{document}
